\section*{Bab \ref{ch:b1:p-block-16-18} --- Blok p II: Oksigen, Halogen, dan Gas Mulia}

\begin{solution}{exo:b1:p-block-16-18:1}
\ce{H2S} $-$II; \ce{S8} 0; \ce{SO2} $+$IV; \ce{SO3^2-} $+$IV; \ce{SO4^2-}
$+$VI; \ce{S2O3^2-} $+$II rata-rata; \ce{H2S2O7} $+$VI.
\end{solution}

\begin{solution}{exo:b1:p-block-16-18:2}
\omterm{def:b1:p-block-16-18:halogen}{Halogen} mengoksidasi halida di bawahnya: \ce{Cl2 + 2I- -> 2Cl- + I2} dan
\ce{Cl2 + 2Br- -> 2Cl- + Br2} terjadi; iodin dengan bromida dan brom
dengan klorida tidak bereaksi.
\end{solution}

\begin{solution}{exo:b1:p-block-16-18:3}
\ce{SO2}: S dengan \omterm{def:b1:lewis-resonance-vsepr:lewis}{sepasang elektron bebas}, satu S=O dan satu \ce{S+-O-}
dalam masing-masing dari dua \omterm{def:b1:lewis-resonance-vsepr:resonance}{struktur resonansi} (atau dua S=O dengan oktet
yang diperluas): bengkok. \ce{SO3}: segitiga datar, tiga ikatan S--O yang
ekuivalen. \ce{O3}: O pusat dengan \omterm{def:b1:lewis-resonance-vsepr:lewis}{sepasang elektron bebas}, satu O=O dan
satu O--O$^-$, dua \omterm{def:b1:lewis-resonance-vsepr:resonance}{struktur resonansi}: bengkok.
\end{solution}

\begin{solution}{exo:b1:p-block-16-18:4}
\ce{SO2 + H2O <=> H2SO3}; \ce{SO2 + 2NaOH -> Na2SO3 + H2O};
\ce{SO3 + H2O -> H2SO4}; \ce{SO3 + 2NaOH -> Na2SO4 + H2O}.
\end{solution}

\begin{solution}{exo:b1:p-block-16-18:5}
HF: $x^2/(0.10 - x) = 10^{-3.16}$, $x = \qty{8.0e-3}{mol/L}$, pH 2.10. HCl:
pH 1.00. Ikatan H--F kuat dan ion fluorida terikat kuat oleh \omterm{def:b1:intermolecular-forces-solvents:hbond}{ikatan
hidrogen}: HF adalah \omterm{def:b1:predominant-reaction:strong-acid}{asam lemah}.
% ledger: dfg:HF_ao, dfg:F-_ao
\end{solution}

\begin{solution}{exo:b1:p-block-16-18:6}
\ce{BrF5}: lima ikatan dan \omterm{def:b1:lewis-resonance-vsepr:lewis}{sepasang elektron bebas}, piramida segi empat.
\ce{IF7}: bipiramida segi lima. \ce{ICl4-}: empat ikatan dan dua \omterm{def:b1:lewis-resonance-vsepr:lewis}{pasangan
elektron bebas}, segi empat datar. \ce{XeO3}: tiga ikatan dan \omterm{def:b1:lewis-resonance-vsepr:lewis}{sepasang
elektron bebas}, piramida segitiga.
\end{solution}

\begin{solution}{exo:b1:p-block-16-18:7}
$\log K = 2(1.36 - 0.53)/0.059 = 28$. Iodin yang dibebaskan membentuk
\omterm{def:b1:complexation:complex}{kompleks} berwarna biru tua dengan kanji.
\end{solution}

\begin{solution}{exo:b1:p-block-16-18:8}
$E^\circ(\ce{O3}/\ce{O2}) = 2.07 > 0.53$~V: \ce{O3 + 2I- + 2H+ -> O2 + I2
+ H2O}. Iodin yang terbentuk dititrasi dengan tiosulfat
(\cref{exo:b1:nernst:7}): satu \ce{I2} per \ce{O3}.
\end{solution}

\begin{solution}{exo:b1:p-block-16-18:9}
\ce{C12H22O11 -> 12C + 11H2O}; $M = \qty{342.0}{g/mol}$, $10.0/342.0 =
\qty{0.0292}{mol}$, karbon $0.0292 \times 12 \times 12.0 = \qty{4.21}{g}$.
\end{solution}

\begin{solution}{exo:b1:p-block-16-18:10}
$x(0.10 + x)/(0.10 - x) = 10^{-1.99}$: $x = \qty{8.6e-3}{mol/L}$,
$[\ce{H3O+}] = 0.1086$, $\mathrm{pH} = 0.96$ dan bukan 1.00: keasaman
kedua menambah sekitar \qty{9}{\percent} pada $[\ce{H3O+}]$.
\end{solution}

\begin{solution}{exo:b1:p-block-16-18:11}
\ce{F2}/\ce{F-} (2.89~V) terletak di atas semua pasangan lain: tidak ada
\omterm{def:b1:nernst:couple}{oksidator} yang dapat mengambil elektron dari fluorida. Pada elektrolisis
larutan berair, air, yang teroksidasi pada potensial yang jauh lebih
rendah (\ce{O2}/\ce{H2O}, 1.23~V), akan bereaksi sebagai gantinya; leburan
fluorida tidak mengandung air.
\end{solution}

\begin{solution}{exo:b1:p-block-16-18:12}
$\Delta_r G^\circ = -51.4 - (16.40 - 51.57) = \qty{-16.2}{kJ/mol}$,
$\log K = 16.2/5.708 = 2.84$, $K = 7 \times 10^2$.
% ledger: dfg:I2_ao, dfg:I-_ao, dfg:I3-_ao
\end{solution}

\begin{solution}{pb:b1:p-block-16-18:1}
\textbf{1.} 0, $+$IV, $+$VI, $+$VI.
\textbf{2.} \ce{S + O2 -> SO2}.
\textbf{3.} S dengan dua domain ikatan dan \omterm{def:b1:lewis-resonance-vsepr:lewis}{sepasang elektron bebas}:
bengkok, sekitar \ang{120}.
\textbf{4.} $10^6/32.1 = \qty{3.12e4}{mol}$.
\textbf{5.} \qty{3.115e4}{mol} \ce{O2}: $V = 3.115 \times 10^4 \times
8.314 \times 298.15/10^5 = \qty{772}{m^3}$, udara $772/0.21 = \qty{3.7e3}{m^3}$.
\textbf{6.} Pembakarannya sangat eksoterm; \omterm{def:b1:elementary-steps:catalyst}{katalis} bekerja dalam rentang
suhu yang terbatas, dan oksidasi \ce{SO2}, yang juga eksoterm, kurang
sempurna dalam keadaan panas.
\textbf{7.} \ce{2SO2 + O2 <=> 2SO3}.
\textbf{8.} Tetapan kesetimbangan yang besar tidak mengatakan apa pun
tentang laju: pada suhu kamar reaksinya tidak berjalan.
\textbf{9.} \omterm{def:b1:elementary-steps:catalyst}{Katalis} itu menyediakan jalan dengan \omterm{def:b1:rate-laws:activation-energy}{energi aktivasi} yang
lebih rendah dan terbentuk kembali: sebuah \omterm{def:b1:elementary-steps:catalyst}{katalis} heterogen.
\textbf{10.} Segitiga datar.
\textbf{11.} Oksigen yang lebih banyak menjaga kuosien
$Q = p_{\ce{SO3}}^2/(p_{\ce{SO2}}^2 p_{\ce{O2}})$ tetap lebih rendah:
konversi \ce{SO2} berlangsung lebih jauh.
\textbf{12.} \qty{0.3}{\percent}.
\textbf{13.} \ce{SO3 + H2SO4 -> H2S2O7}.
\textbf{14.} Trioksida itu membentuk kabut tetesan asam yang halus yang lolos
melewati penyerap.
\textbf{15.} \ce{H2S2O7 + H2O -> 2H2SO4}.
\textbf{16.} \ce{2S + 3O2 + 2H2O -> 2H2SO4}.
\textbf{17.} Satu molekul air per atom belerang: $\num{3.115e4} \times 18.0 =
\qty{561}{kg}$.
\textbf{18.} Pengenceran melepaskan banyak kalor; air yang dituangkan ke
atas asam langsung mendidih dan memercikkan asam.
\textbf{19.} $x(0.050 + x)/(0.050 - x) = 10^{-1.99}$: $x = \qty{7.5e-3}{mol/L}$,
$[\ce{H3O+}] = \qty{0.0575}{mol/L}$, $\mathrm{pH} = 1.24$.
\textbf{20.} $3.12 \times 10^4 \times 0.995 \times 98.1 = \qty{3.04}{t}$.
\textbf{21.} $84 \times 98.1/32.1 = \qty{257}{Mt}$.
\textbf{22.} $98.1/32.1 =$ \textbf{\qty{3.06}{t} asam sulfat per ton
belerang}.
\end{solution}
